\markedsection{\toctheory}{Sequential ACE construction}
\label{app:sequential-ace}

The ACE builder introduced in Sec.~\ref{subsec:pt-builders} constructs the environmental influence by adding one microscopic mode at a time and compressing the resulting temporal bonds.
Here we describe this update in more detail, following the original ACE construction~\cite{Cygorek2022} and its implementation in the ACE toolkit~\cite{CygorekGauger2024}, 
The system operations remain open during construction, so the resulting PT-MPO can subsequently be used with different instrument sequences. 

Figure~\ref{fig:sequential-ace} shows the join for four modes and three propagation intervals, with time running from right to left. 
Each row is one mode, shaded from dark (first bath mode) to light (last bath mode). 
We assume an initially factorised state $\rho_S(0)\otimes\bigotimes_m\rho_{E_m}(0)$ and the generator decomposition given in Sec.~\ref{subsec:pt-builders}: each $\mathcal L_m\equiv\mathcal L_{E_m+SE_m}$ acts on the system and mode $m$, with no direct interactions between modes. 
Joining successive bath indices, contracting the first with $|\rho_{E_m}(0)\rangle\rangle$, and closing the last with $\langle\langle I_{E_m}|$ gives a single-mode influence chain. 
The first step absorbs \(E_1\) and \(E_2\) into a single chain; the same join is then repeated for \(E_3\) and \(E_4\). 
Its internal bath index has dimensions $d_m^2$, where $d_m$ is the mode Hilbert-space dimension. 
As we compress the modes, we lose the "double-wire" space and replace it with a compressed influence, with memory bond links taking its place. 
After every mode has been absorbed and the temporal bonds compressed, the network is the PT-MPO with cores \(Q^{[k]}\). 

%
\begin{figure}[tbp]
    \centering
    \resizebox{\linewidth}{!}{% Sequential ACE construction.  Time runs right to left.
% (a) independent mode chains, (b) partial join, (c) compressed PT-MPO.
% System legs stay open.  Bath lines are Liouville wires; left ends are trace-outs.
% Requires pt_tikz_styles.tex in the preamble.
\begin{tikzpicture}[PTPicture]
  \def\ACEBoxW{5.8mm}%
  \def\ACEBoxH{5.8mm}%
  \def\ACEdx{1.70}%
  \def\ACEdy{0.86}%
  % One red, light for the last mode (top) and dark for the first (bottom).
  \definecolor{ACERed}{HTML}{E07A4C}
  \colorlet{ACEMode3}{ACERed!28}
  \colorlet{ACEMode2}{ACERed!46}
  \colorlet{ACEMode1}{ACERed!68}
  \colorlet{ACEMode0}{ACERed}
  \tikzset{
    aceidx/.style={
      font=\scriptsize,
      text=PTLabel,
      inner sep=0.4pt,
      anchor=north,
      yshift=-1.35mm,
    },
    aceenv/.style={
      font=\small,
      text=PTLabel,
      inner sep=0pt,
      anchor=south west,
      xshift=-0.55mm,
      yshift=-0.35mm,
    },
    acearrow/.style={
      draw=PTLabel,
      line width=0.7pt,
      -{Stealth[length=2.1mm, width=1.5mm]},
    },
    acebond/.style={
      draw=PTBond,
      line width=1.7pt,
      line cap=butt,
      on background layer,
    },
    acemem/.style={
      draw=PTBond,
      line width=\PTLinkLW,
      line cap=butt,
      line join=round,
      on background layer,
    },
  }

  % Three propagators.  k=0 is the earliest column (right); k=2 is the latest (left).
  \newcommand{\ACERow}[3]{%
    \foreach \k [evaluate=\k as \xx using (2-\k)*\ACEdx] in {0,1,2} {%
      \PTTensor[
        fill=#3,
        minimum width=\ACEBoxW,
        minimum height=\ACEBoxH,
        font=\scriptsize,
      ]{#1-\k}{(\xx,#2)}{}%
    }%
    \draw[PTLiouvilleLeg] (#1-2.east) -- (#1-1.west);
    \draw[PTLiouvilleLeg] (#1-1.east) -- (#1-0.west);
  }
  % Bath boundaries: short trace-out on the left, labelled initial state on the right.
  % #2 is the mode index printed as E_{#2}.
  \newcommand{\ACEBathEnds}[2]{%
    \PTTraceOut{#1-tr}{($ (#1-2.west)+(-2.7mm,0) $)}%
    \draw[PTLiouvilleLeg] (#1-2.west) -- (#1-tr-east);
    \node[
      name=#1-init,
      inner sep=0pt,
      outer sep=0pt,
      minimum width=3.5mm,
      minimum height=4.8mm,
      anchor=west,
    ] at ($ (#1-0.east)+(2.2mm,0) $) {};
    \draw[PTLiouvilleLeg] (#1-0.east) -- (#1-init.west);
    \begin{scope}[on PTforeground layer]
      \PTSetTriangleRound{3.5mm}{4.8mm}%
      \filldraw[
        draw=black,
        fill=PTLiouville!16,
        line width=\PTCoreStroke,
        line join=round,
        rounded corners=\PTRelRound,
      ] (#1-init.west |- #1-init.south) -- (#1-init.east)
        -- (#1-init.west |- #1-init.north) -- cycle;
      \node[aceenv] at (#1-init.north east) {\(E_{#2}\)};
    \end{scope}
  }
  \newcommand{\ACEVert}[2]{%
    \foreach \k in {0,1,2} {%
      \draw[PTLiouvilleLeg] (#1-\k.south) -- (#2-\k.north);
    }%
  }
  % Open system pair on a stacked column. Both open ends sit below the
  % first (bottom) mode. The input leaves its south edge and bends right;
  % the output leaves the last (top) mode and is brought down the left side.
  \newcommand{\ACEStackLegs}[4]{%
    \coordinate (#1-rail) at ([yshift=-2.8mm]#2.south);
    \coordinate (#1-idrop) at (#2.south);
    \coordinate (#1-ibend) at (#1-idrop |- #1-rail);
    \coordinate (#1-istub) at ($(#1-ibend)+(1.2mm,0)$);
    \coordinate (#1-iend) at ($(#1-istub)+(2.75mm,0)$);
    \draw[PTLiouvilleLeg, rounded corners=1.1mm]
      (#1-idrop) -- (#1-ibend) -- (#1-istub);
    \draw[PTLiouvilleLeg primed] (#1-istub) -- (#1-iend);
    \coordinate (#1-rise) at ($(#3.north)+(0,1.5mm)$);
    \coordinate (#1-side) at ($(#3.north west)+(-1.7mm,1.5mm)$);
    \coordinate (#1-low) at (#1-side |- #1-rail);
    \coordinate (#1-oend) at ($(#1-low)+(-1.7mm,0)$);
    \draw[
      PTLiouvilleLeg,
      rounded corners=1.1mm,
      preaction={draw=white, line width=\PTBridgeHalo, double=none},
    ] (#3.north) -- (#1-rise) -- (#1-side) -- (#1-low) -- (#1-oend);
    \node[aceidx] at (#1-oend) {\(o_{#4}\)};
    \node[aceidx] at (#1-iend) {\(i_{#4}\)};
  }
  % Both legs leave the bottom of one PT-MPO core and bend outward.
  \newcommand{\ACEOpenLegs}[2]{%
    \coordinate (#1-rail) at ([yshift=-2.8mm]#1.south);
    \coordinate (#1-odrop) at ([xshift=-1.35mm]#1.south);
    \coordinate (#1-idrop) at ([xshift=1.35mm]#1.south);
    \coordinate (#1-obend) at (#1-odrop |- #1-rail);
    \coordinate (#1-ibend) at (#1-idrop |- #1-rail);
    \coordinate (#1-ostub) at ($(#1-obend)+(-1.2mm,0)$);
    \coordinate (#1-istub) at ($(#1-ibend)+(1.2mm,0)$);
    \coordinate (#1-oend) at ($(#1-ostub)+(-1.5mm,0)$);
    \coordinate (#1-iend) at ($(#1-istub)+(1.5mm,0)$);
    \draw[PTLiouvilleLeg, rounded corners=1.1mm]
      (#1-odrop) -- (#1-obend) -- (#1-ostub);
    \draw[PTLiouvilleLeg] (#1-ostub) -- (#1-oend);
    \draw[PTLiouvilleLeg, rounded corners=1.1mm]
      (#1-idrop) -- (#1-ibend) -- (#1-istub);
    \draw[PTLiouvilleLeg primed] (#1-istub) -- (#1-iend);
    \node[aceidx] at (#1-oend) {\(o_{#2}\)};
    \node[aceidx] at (#1-iend) {\(i_{#2}\)};
  }

  % ===================================================================== (a)
  \begin{scope}[shift={(0.15,0.55)}]
    \ACERow{a3}{3*\ACEdy}{ACEMode3}
    \ACERow{a2}{2*\ACEdy}{ACEMode2}
    \ACERow{a1}{\ACEdy}{ACEMode1}
    \ACERow{a0}{0}{ACEMode0}
    \ACEBathEnds{a3}{4}
    \ACEBathEnds{a2}{3}
    \ACEBathEnds{a1}{2}
    \ACEBathEnds{a0}{1}
    \ACEVert{a3}{a2}
    \ACEVert{a2}{a1}
    \ACEVert{a1}{a0}
    \foreach \k in {0,1,2} {%
      \ACEStackLegs{a-\k}{a0-\k}{a3-\k}{\k}%
    }
  \end{scope}

  % ===================================================================== (b)
  \begin{scope}[shift={(6.91,0.55)}]
    \ACERow{b2}{2*\ACEdy}{ACEMode3}
    \ACERow{b1}{\ACEdy}{ACEMode2}
    \ACEBathEnds{b2}{4}
    \ACEBathEnds{b1}{3}
    \ACERow{bj}{0}{ACEMode1}
    \ACEVert{b2}{b1}
    \ACEVert{b1}{bj}
    \draw[acebond] (bj-2.east) -- (bj-1.west);
    \draw[acebond] (bj-1.east) -- (bj-0.west);
    \foreach \k in {0,1,2} {%
      \ACEStackLegs{b-\k}{bj-\k}{b2-\k}{\k}%
    }
  \end{scope}

  % ===================================================================== (c)
  \begin{scope}[shift={(14.00,0.78)}]
    \foreach \k [evaluate=\k as \xx using (2-\k)*\ACEdx] in {0,1,2} {%
      \PTTensor[
        fill=QAMPOcore,
        minimum width=7.6mm,
        minimum height=12.4mm,
        font=\scriptsize,
      ]{c-\k}{(\xx,0.22)}{\(Q^{[\k]}\)}%
    }%
    \draw[acemem] ([yshift=2.4mm]c-2.east) -- ([yshift=2.4mm]c-1.west);
    \draw[acemem] ([yshift=2.4mm]c-1.east) -- ([yshift=2.4mm]c-0.west);
    \foreach \k in {0,1,2} {%
      \ACEOpenLegs{c-\k}{\k}%
    }
  \end{scope}

  % Each arrow sits at the vertical midpoint of the panel it points to.
  \draw[acearrow] (5.04,1.41) -- (5.87,1.41);
  \draw[acearrow] (11.80,1.00) -- (13.20,1.00);
  \node[font=\small, text=PTLabel, inner sep=0.5pt, anchor=south]
    at (12.50,1.04) {Repeat};
  \node[font=\small, text=PTLabel, inner sep=0.5pt, anchor=north]
    at (12.50,0.96) {\(E_3,E_4\)};
\end{tikzpicture}
}
    \caption{Sequential construction of the PT-MPO, illustrated for four bath modes and three propagation intervals, with time running from right to left. 
    Each bath mode $E_m$ contains an initial state and its propagator across time. 
    The vertical legs of each bath mode contain the influence of that bath mode on the system. 
    By merging this leg sequentially for each bath mode with SVD truncation, we can retain only the relevant influence on the system at each timestep. 
    Repeating this over all bath modes gives the final PT-MPO. 
    Figure adapted from Lacroix \textit{et al.}~\cite{LacroixReview2026}.}
    \label{fig:sequential-ace}
\end{figure}
%

Let $\mathcal L_{<m}=\sum_{j<m}\mathcal L_j$ denote the contribution of the modes already included. 
The asymmetric addition of mode $m$ uses
%
\begin{equation}
    e^{\Delta t(\mathcal L_{<m}+\mathcal L_m)}
    = e^{\Delta t\mathcal L_m}
      e^{\Delta t\mathcal L_{<m}}
      +\mathcal O(\Delta t^2).
    \label{eq:ace-asymmetric-split}
\end{equation}
%
The symmetric addition of mode $m$ uses
%
\begin{equation}
    e^{\Delta t(\mathcal L_{<m}+\mathcal L_m)}
    = e^{\frac{\Delta t}{2}\mathcal L_m}
      e^{\Delta t\mathcal L_{<m}}
      e^{\frac{\Delta t}{2}\mathcal L_m}
      +\mathcal O(\Delta t^3).
    \label{eq:ace-symmetric-split}
\end{equation}
%
This decomposition does not matter if we have commuting interaction terms of each bath mode with the system. 
However, if the bath modes are non-commutative in nature, the latter would lead to a more accurate compression. 
Compression reduces the enlarged temporal bonds while we merge the two MPO chains. 
At a sweep step, the appropriate core and any factor carried from its neighbour are reshaped into a matrix and decomposed as
%
\begin{equation}
    M=U\Sigma V^\dagger
    \;\longrightarrow\;
    U_r\Sigma_rV_r^\dagger,
    \qquad
    \sigma_j>\epsilon\sigma_1,\qquad r\leq\chi_{\max}.
    \label{eq:ace-appendix-svd}
\end{equation}
%
One factor becomes the updated core, while the remaining factor is absorbed into the next core in the sweep direction. 
Sweeping through the chain carries information from the initial state and final trace into the intermediate tensors~\cite{Cygorek2022}. 
Its dimension is determined by how well that influence can be compressed over the chosen time interval. 

The two schedules compared in Sec.~\ref{subsec:ace-compressors} differ in when this truncation is performed. The \texttt{:canonzip} schedule completes the mode join without truncation, canonicalises towards the final time, and then performs a reverse relative-SVD sweep. 
The \texttt{:zipup\_cpp} schedule truncates during the forward join and subsequently sweeps backward, reducing intermediate storage but choosing truncation bases before the complete joined chain has been canonicalised. 
The same numerical cutoff therefore need not give the same accuracy in both schedules. 
We refer to Sec.~\ref{subsec:ace-compressors} for their measured accuracy and computational cost. 

Driven or dissipative modes modify the first ingredient, while tree-like ACE joins partial influences hierarchically and introduces preselection and further compression sweeps~\cite{Cygorek2024Treelike}.
These are the extension points described in Sec.~\ref{sec:roadmap} and can retain the same final \texttt{ProcessTensor} representation.
Initially correlated modes or direct intermode interactions instead require revisiting the independent-mode decomposition used above.
